\scanpage{138}
\[
\begin{aligned}
&\leq\left(\sum_n\alpha_n(H^m)\rho^{np'r}\right)
\left(\sum_n\alpha_n(K^m)\rho^{nq'r}\right)\\
&\deleted{\unsure{=|\leq}}\qquad=M_{H^m}(\rho^{p'r})M_H(\rho^{q'r})
\end{aligned}
\]

d'où, compte tenu de (12)

\[
\frac{\lambda(r/2m)}{r/2m}S_r(AB)\leq
\int_0^\infty\frac{\log M_{H^m}(\rho^{p'r})}{\rho^{1+r/2m}}\,d\rho+
\int_0^\infty\frac{\log M_{K^m}(\rho^{q'r})}{\rho^{1+r/2m}}\,d\unsure{\rho}
\tag*{(13)}
\]

Calculons \deleted{\unsure{par exemple}} la première intégrale du second membre, par le changement de variable $\sigma=\rho^{p'r}$ d'où $\rho=\sigma^{pr'}$ $\frac{d\rho}{\rho}=pr'\frac{d\sigma}{\sigma}$ il vient

\[
\int_0^\infty\frac{\log M_{H^m}(\rho^{p'r})}{\rho^{r/2m}}\frac{d\rho}{\rho}
=pr'\int_0^\infty\frac{\log M_{H^m}(\sigma)}{\sigma^{p/2m}}\frac{d\sigma}{\sigma}
\]

Le deuxième membre est égal (\unsure{\deleted{formule}|formule} (9)) à

\[
pr'\frac{\lambda(p/2m)}{p/2m}S_{p/2m}(H^m)=\frac{\lambda(\unsure{p|p/2m})}{r/2m}S_p(A).
\]

\unsure{De} $=$ la deuxième intégrale dans (13) est égale à $\frac{\lambda(q/2m)}{r/2m}S_q(B)$, et \deleted{\unsure{on obtient|(13) est}} \inserted{(13) s'écrit \unsure{ainsi}}

\[
\lambda(r/2m)S_r(AB)\leq\lambda(p/2m)S_p(A)+\lambda(q/2m)S_q(B)
\]

ou aussi

\[
S_r(AB)\leq\frac{\sin\pi\frac{r}{2m}}{\sin\pi\frac{p}{2m}}S_p(A)+
\frac{\sin\pi\frac{r}{2m}}{\sin\pi\frac{q}{2m}}S_q(B)
\]

\inserted{Cela prouve déjà que $S_{\unsure{r|p}}(AB)\leq+\infty$, et donne même une majoration effective de $S_r(AB)$ en fonction de $S_p(A)$ et $S_q(B)$.}

Faisant tendre $m$ vers $+\infty$, on obtient la formule plus simple

\[
\boxed{S_r(AB)\leq\frac{r}{\deleted{2m}p}S_p(A)+\frac rq S_q(B)}\tag*{(14)}
\]

C'est une inégalité remarquable, mais n'est pas celle que nous cherchons. Mais \deleted{\unsure{cette relation} \illegible{1}}\textsuperscript{\deleted{\illegible{1}}}
